% Part: normal-modal-logic
% Chapter: filtrations
% Section: S5-decidable

\documentclass[../../../include/open-logic-section]{subfiles}

\begin{document}

\olfileid{nml}{fil}{dec}

\olsection{المنطق \Log{S5} قابل للقرار}
بخاصية النموذج المنتهي، إذا اقترنت بتعداد فعال للبراهين وبامتحان فعال
للنماذج المنتهية المقصودة، يتيسر إثبات أن نظمًا من المنطق الجهوي
\emph{قابلة للقرار}. ومعنى ذلك وجود إجراء قابل للحساب يبتّ فيما إذا
كانت !!a{formula} معطاة !!{derivable} في النظام أم لا.
% تصحيح دلالي (0457-I1): الخاصية وحدها لا تكفي؛ يلزم أن تكون جهتا البحث قابلتين للتنفيذ الفعلي.

\begin{thm}
  \Log{S5} قابل للقرار.
\end{thm}

\begin{proof}
  لنأخذ~$!A$، ولتكن ال!!{propositional variable}s التي ترد في
  $!A$ واقعة بين~$p_1$، \dots،~$p_k$. لكل عدد~$n$ عدد منتهٍ من
  النماذج الكلية ذات~$n$ من العوالم ومن تقويمات~$p_1$،
  \dots،~$p_k$ عليها. فنُجري تعدادين \emph{بالتوازي}: نولّد البراهين
  على ترتيب منتظم، فنعدّ بذلك مبرهنات~\Log{S5} كلها؛ ونعدّ النماذج
  الكلية المنتهية، مبتدئين بذوات~$1$ من العوالم، ثم~$2$ ثم~$\dots$،
  ونختبر في كل منها هل تكذب~$!A$ عند عالم من عوالمه.
  ولا بد أن ينتهي أحد البحثين إلى الجواب. فإن
  \olref[com][fra]{thm:generaldet} و\olref[fmp]{cor:S5fmp}
  تقتضيان أن~$!A$ إما !!{derivable} وإما كاذبة في نموذج كلي منتهٍ.
\end{proof}

وإنما استقام هذا البرهان لـ~\Log{S5} لأن ترشيح النموذج الكلي كلي
بالضرورة. وكذلك تُحفظ الانعكاسية والتسلسل بالترشيح، وأما سائر الخواص
فتحتاج إلى نظر زائد.

\end{document}
